% CSMATH 2026 Course Paper Template
% 浙江大学计算机科学与技术学院 · 计算机应用数学
% School of Computer Science & Technology, Zhejiang University

\documentclass[11pt,a4paper,UTF8]{ctexart}

% === 页边距设置 (A4) ===
\usepackage[top=2.5cm, bottom=2.5cm, left=2.5cm, right=2.5cm]{geometry}

% === 导入宏包 ===
% ============================================================
% CSMATH 2026 课程论文 · 导言区
% 浙江大学计算机科学与技术学院 · 计算机应用数学
% ============================================================

% --- 基础宏包 ---
\usepackage{xcolor}

% --- 数学宏包 ---
\usepackage{amsmath,amssymb,amsthm}
\usepackage{mathtools}

% --- 图表与图形 ---
\usepackage{graphicx}
\usepackage{booktabs}
\usepackage{caption}
\usepackage{subcaption}

% --- 算法 ---
\usepackage[ruled,linesnumbered]{algorithm2e}
\SetKwComment{Comment}{$\triangleright$\ }{}

% --- 定理环境 ---
\newtheorem{theorem}{定理}[section]
\newtheorem{lemma}[theorem]{引理}
\newtheorem{proposition}[theorem]{命题}
\newtheorem{corollary}[theorem]{推论}
\newtheorem{definition}{定义}[section]
\newtheorem{remark}{注记}[section]
\newtheorem{property}{性质}[section]
\theoremstyle{definition}
\newtheorem{example}{例}[section]

% --- 常用数学运算符 ---
\DeclareMathOperator*{\argmin}{arg\,min}
\DeclareMathOperator*{\argmax}{arg\,max}
\DeclareMathOperator*{\E}{\mathbb{E}}
\DeclareMathOperator*{\Var}{Var}
\DeclareMathOperator*{\Cov}{Cov}
\DeclareMathOperator{\diag}{diag}
\DeclareMathOperator{\tr}{tr}
\DeclareMathOperator{\rank}{rank}

% --- 常用符号 ---
\newcommand{\R}{\mathbb{R}}
\newcommand{\C}{\mathbb{C}}
\newcommand{\N}{\mathbb{N}}
\newcommand{\Z}{\mathbb{Z}}
\newcommand{\X}{\mathcal{X}}
\newcommand{\Y}{\mathcal{Y}}
\newcommand{\F}{\mathcal{F}}
\newcommand{\Loss}{\mathcal{L}}
\newcommand{\norm}[1]{\left\lVert#1\right\rVert}
\newcommand{\inner}[2]{\left\langle#1, #2\right\rangle}
\newcommand{\grad}{\nabla}

% --- 行内标注（写作过程中使用，提交前禁用）---
\newcommand{\red}[1]{{\color{red}#1}}
\newcommand{\todo}[1]{{\color{red}#1}}
\newcommand{\TODO}[1]{\textbf{\color{red}[TODO: #1]}}
% 提交版本取消注释下面两行以隐藏 TODO
% \renewcommand{\TODO}[1]{}
% \renewcommand{\todo}[1]{#1}

% --- 微排版优化（定稿前取消注释）---
% \usepackage{microtype}

% --- 其他 ---
\usepackage{enumitem}
\usepackage{multirow}
\usepackage{url}


% === 超链接设置 (模板要求使用 hyperref) ===
\definecolor{cvprblue}{rgb}{0.21,0.49,0.74}
\usepackage[pagebackref,breaklinks,colorlinks,allcolors=cvprblue]{hyperref}

% ==================== 论文信息（请填写）====================
\newcommand{\studentName}{姓名}
\newcommand{\studentId}{学号}
\newcommand{\courseName}{计算机应用数学}
\newcommand{\courseCode}{CSMATH 2026}
\newcommand{\semester}{2026 春夏学期}

% ==================== 标题 ====================
\title{在此处填写论文标题\\[4pt]
\large \textnormal{CSMATH 2026 课程论文}}

% ==================== 作者 ====================
\author{\studentName\\[4pt]
{\small 学号：\studentId}\\[4pt]
{\small \courseName\ (\courseCode)}\\[4pt]
{\small 浙江大学计算机科学与技术学院}\\[4pt]
{\small \semester}
}

\date{}

\begin{document}

\maketitle

% ==================== 摘要 ====================
% ============================================================
% 摘要
% ============================================================
\begin{abstract}
在此处撰写摘要（中文约 200--400 字，或英文约 150--250 词）。
摘要应简明扼要地涵盖以下要点：
\begin{enumerate}[nosep, leftmargin=*]
    \item 研究问题是什么？（一句话点明应用背景与数学问题类型）
    \item 本文采用的核心数学工具/方法是什么？（与课程内容的关联）
    \item 主要结果或发现是什么？（包括理论分析与实验结果）
\end{enumerate}

\noindent\textbf{关键词：}计算机应用数学；【补充 3--5 个关键词】
\end{abstract}


% ==================== 正文 ====================
% ============================================================
% 一、引言 (Introduction)
% ============================================================
\section{引言}

\subsection{研究背景与动机}

在此处简要介绍应用背景：这个问题的实际来源（可以是图形学、视觉、机器学习、数据科学等领域的真实场景），为什么它是一个值得研究的数学问题。

\begin{itemize}
    \item 描述问题的应用场景和实际意义；
    \item 指出现有方法的局限性或尚未解决的理论问题；
    \item 明确本文要探讨的具体问题范围。
\end{itemize}

\subsection{相关工作}

简要回顾与该问题相关的已有工作，重点关注其中与本课程数学主题相关的部分。

\begin{itemize}
    \item 文献综述应聚焦于数学方法，而非仅是实验结果对比；
    \item 明确指出已有工作与本课程哪些章节相关；
    \item 引用应优先来自顶会/顶刊或 arXiv。
\end{itemize}

% 示例引用格式（删除或替换）：
% 相关方法最早由 \cite{Alpher02} 提出，后续工作 \cite{Alpher03,Alpher04} 从优化角度进行了改进。本课程讲义第 X 章详细讨论了该方法的核心思想。

\subsection{本文贡献}

用条目清晰列出本文的主要工作（3--4 点即可）：

\begin{enumerate}
    \item 对【具体问题】给出了严格的形式化定义，明确了变量、假设、目标与约束；
    \item 从【几何/物理/概率】视角给出了核心数学对象的直观解释；
    \item 推导了【具体算法或分析方法】，分析了其收敛性/复杂度/数值稳定性；
    \item 实现了核心算法并设计了可复现实验，验证了方法的有效性及边界条件。
\end{enumerate}

\subsection{论文结构}

简述后续各节的安排，与课程要求的能力主线（形式化→直观化→数学化→可计算）相呼应。
           % 一、引言
% ============================================================
% 二、问题形式化 (Problem Formalization)
% ============================================================
\section{问题形式化}

\subsection{变量与符号定义}

用严格的数学符号定义问题中涉及的所有变量。建议使用表格整理：

\begin{table}[htbp]
    \centering
    \caption{符号定义表}
    \label{tab:notation}
    \begin{tabular}{c|l}
        \toprule
        符号 & 含义 \\
        \midrule
        $\X \subseteq \R^d$ & 输入/数据空间 \\
        $\Y$                & 输出/标签空间 \\
        $\mathbf{x}_i$      & 第 $i$ 个样本 \\
        $\boldsymbol{\theta}$ & 待优化参数 \\
        $\Loss(\cdot)$      & 损失/目标函数 \\
        \bottomrule
    \end{tabular}
\end{table}

\subsection{假设}

明确列出为简化问题或使推导成立所引入的所有假设，每一条都要有明确的数学表述：

\begin{enumerate}[label=(\textbf{A\arabic*}), leftmargin=*]
    \item \textbf{数据假设}：【例如：样本独立同分布、数据满足某种光滑性条件等】
    \item \textbf{模型假设}：【例如：线性模型、高斯噪声、低秩结构等】
    \item \textbf{优化假设}：【例如：目标函数 $L$-光滑、$\mu$-强凸、约束可行域非空等】
\end{enumerate}

\subsection{目标函数与约束}

写出完整的优化问题或数学模型的数学表达式：

\begin{equation}\label{eq:problem}
    \min_{\mathbf{x} \in \X} \; f(\mathbf{x}) \quad \text{s.t.} \quad g_i(\mathbf{x}) \leq 0,\; i = 1,\ldots,m
\end{equation}

或对于统计推断/学习问题，写出概率模型和推断目标。

\subsection{问题类别}

指出问题所属的数学类别（勾选适用的项目）：

\begin{itemize}
    \item 凸 / 非凸
    \item 连续 / 离散 / 混合整数
    \item 确定 / 随机
    \item 约束 / 无约束
    \item 光滑 / 非光滑
    \item 凸优化 / 组合优化 / 变分推断 / 谱方法 / PDE / 泛函分析
\end{itemize}

这一判断直接影响后续方法的选择与分析路径。
   % 二、问题形式化
% ============================================================
% 三、数学背景与直觉 (Math Background & Intuition)
% ============================================================
\section{数学背景与直觉}

% 提示：本节的核心目标是"用一句话、一张图说明它究竟在做什么"，
% 而不是堆砌公式。多用自己的话和具体的小例子。

\subsection{核心概念概述}

列出本文涉及的核心数学概念，每个概念用一段话说明它"是什么"以及"为什么在此处有用"。

\begin{itemize}
    \item \textbf{【概念 1，例如：梯度下降】}：【用自己的话解释梯度下降的几何意义：每次沿局部最陡下降方向移动，就像蒙着眼睛下山时每次朝脚底最陡的方向迈一步。关联课程第 X 章。】
    \item \textbf{【概念 2，例如：对偶问题】}：【用自己的话解释对偶的物理/几何直觉……关联课程第 Y 章。】
    \item \textbf{【概念 3，…】}：【……】
\end{itemize}

\subsection{直观解释}

\subsubsection{一个小例子}

用一个极度简化的小规模例子（2--3 维，手算级别）展示核心方法到底在做什么。

\begin{example}
    考虑一个二维问题：
    \begin{equation}
        \min_{x_1, x_2} \; f(x_1, x_2) = (x_1 - 1)^2 + 10(x_2 - 2)^2
    \end{equation}
    这是一个典型的病态二次函数……【解释梯度下降在此处的行为：zigzag 现象及其几何原因】
\end{example}

\subsubsection{几何/物理图像}

% === 建议在此处插入一张手绘示意图或矢量图 ===

\begin{figure}[htbp]
    \centering
    % \includegraphics[width=0.6\textwidth]{figures/intuition.pdf}
    \caption{【图的标题：例如，梯度下降在病态二次函数上的 zigzag 轨迹（等高线图）】}
    \label{fig:intuition}
\end{figure}

图~\ref{fig:intuition} 展示了……【解释图中各元素的含义】

\subsection{与课程内容的关联}

\begin{itemize}
    \item 本文的核心数学工具对应课程讲义第【X】章（【具体主题】）；
    \item 此外，推导中用到第【Y】章的【具体知识点】；
    \item 【如果有】课程中未涵盖但本文用到的额外数学知识，简要说明并给出参考文献。
\end{itemize}
       % 三、数学背景与直觉
% ============================================================
% 四、方法分析 (Method Analysis)
% ============================================================
\section{方法分析}

\subsection{算法描述}

给出完整的伪代码或算法步骤：

\begin{algorithm}[htbp]
    \caption{【算法名称】}
    \label{alg:main}
    \SetKwInOut{Input}{输入}
    \SetKwInOut{Output}{输出}
    \Input{【输入变量说明】}
    \Output{【输出变量说明】}
    \BlankLine
    \textbf{初始化}：设置参数 $\boldsymbol{\theta}^{(0)}$，学习率 $\eta$，最大迭代次数 $T$\;
    \For{$t = 0, 1, \ldots, T-1$}{
        计算梯度 $\mathbf{g}^{(t)} = \grad f(\boldsymbol{\theta}^{(t)})$\;
        更新参数 $\boldsymbol{\theta}^{(t+1)} = \boldsymbol{\theta}^{(t)} - \eta \mathbf{g}^{(t)}$\;
        \If{满足收敛条件}{
            终止循环\;
        }
    }
    \Return $\boldsymbol{\theta}^{(T)}$\;
\end{algorithm}

\subsection{推导过程}

给出关键公式的推导。保持步骤清晰，注明每一步的依据（假设、性质或引用的定理）。

\begin{proof}[推导]
    展开推导过程：
    \begin{align}
        \norm{\grad f(\mathbf{x}) - \grad f(\mathbf{y})}
        &\leq L \norm{\mathbf{x} - \mathbf{y}} \quad \text{(由 $L$-光滑假设)} \label{eq:smooth} \\
        f(\mathbf{y}) &\leq f(\mathbf{x}) + \inner{\grad f(\mathbf{x})}{\mathbf{y} - \mathbf{x}} + \frac{L}{2}\norm{\mathbf{y} - \mathbf{x}}^2 \label{eq:descent}
    \end{align}
\end{proof}

\subsection{收敛性/误差分析}

\begin{theorem}[收敛性]\label{thm:convergence}
    在假设 A1--A3 下，算法~\ref{alg:main} 产生的序列 $\{\mathbf{x}^{(t)}\}$ 满足：
    \begin{equation}
        f(\mathbf{x}^{(t)}) - f(\mathbf{x}^*) \leq \frac{L \norm{\mathbf{x}^{(0)} - \mathbf{x}^*}^2}{2t}
    \end{equation}
    即达到 $\epsilon$-精度需要 $O(1/\epsilon)$ 次迭代。
\end{theorem}

\begin{proof}
    由式~(\ref{eq:descent}) 和梯度下降的更新规则……（给出简洁证明）
\end{proof}

\subsection{复杂度与数值稳定性讨论}

\begin{itemize}
    \item \textbf{时间复杂度}：每次迭代 $O(\cdots)$，总计 $O(\cdots)$；
    \item \textbf{空间复杂度}：$O(\cdots)$；
    \item \textbf{数值稳定性}：【讨论关键步骤的数值问题，如矩阵求逆的 conditioning、梯度消失/爆炸等】；
    \item \textbf{参数敏感性}：【讨论学习率等超参数对收敛的影响】。
\end{itemize}
          % 四、方法分析
% ============================================================
% 五、实现与实验 (Implementation & Experiments)
% ============================================================
\section{实现与实验}

\subsection{实现环境}

\begin{itemize}
    \item \textbf{编程语言}：Python 3.x（推荐）/ C++ / Julia / MATLAB
    \item \textbf{关键依赖}：NumPy, SciPy, Matplotlib, PyTorch（如适用）
    \item \textbf{运行环境}：macOS / Linux，CPU / GPU（说明型号）
    \item \textbf{代码仓库}：【GitHub 链接，建议公开仓库并附 README】
\end{itemize}

\subsection{实验设置}

\subsubsection{数据集}

描述所使用的数据集（真实数据或合成数据），包括规模、维度、生成方式。

\subsubsection{评估指标}

列出并解释用于衡量方法性能的定量指标（例如：目标函数值 vs 迭代次数、重构误差、分类准确率、运行时间等）。

\subsubsection{对比方法}

如果进行了对比实验，列出所有对比方法及其简要说明。

\subsection{实验结果}

% === 至少一张图 ===

\begin{figure}[htbp]
    \centering
    % \includegraphics[width=0.75\textwidth]{figures/convergence.pdf}
    \caption{【图 1 标题：例如，目标函数值随迭代次数的收敛曲线】}
    \label{fig:convergence}
\end{figure}

说明：图~\ref{fig:convergence} 展示了……

% === 至少一张表 ===

\begin{table}[htbp]
    \centering
    \caption{【表 1 标题：例如，不同方法的性能对比】}
    \label{tab:results}
    \begin{tabular}{lccc}
        \toprule
        方法 & 指标 1 & 指标 2 & 运行时间 (s) \\
        \midrule
        本文方法 & --- & --- & --- \\
        对比方法 A & --- & --- & --- \\
        对比方法 B & --- & --- & --- \\
        \bottomrule
    \end{tabular}
\end{table}

说明：表~\ref{tab:results} 显示了……

\subsection{结果分析与讨论}

\begin{itemize}
    \item 方法是否按理论预期工作？如不符合，可能的原因是什么？
    \item 参数变化（学习率、正则化系数等）对结果有什么影响？
    \item 方法的成功案例与失败案例分别是什么样的？（建议展示具体例子）
    \item 是否存在理论未覆盖的边界情况？
\end{itemize}
     % 五、实现与实验
% ============================================================
% 六、个人见解 (Personal Insights)
% ============================================================
\section{个人见解}

\subsection{方法总结}

简要总结本文方法的核心思想与关键步骤（用自己的话概括，5--8 行即可）。

\subsection{优点与适用场景}

\begin{itemize}
    \item \textbf{理论上的优势}：【例如：收敛性有严格保证、假设条件宽松等】
    \item \textbf{实践上的优势}：【例如：实现简单、参数少、调参友好等】
    \item \textbf{特别适用的场景}：【描述方法在何种条件下表现最好】
\end{itemize}

\subsection{局限性与改进方向}

\begin{itemize}
    \item \textbf{理论局限}：【例如：强凸假设在实际中不成立时方法可能退化】
    \item \textbf{实践局限}：【例如：高维情况下计算代价过高、对初始化敏感等】
    \item \textbf{可能的改进方向}：
    \begin{enumerate}
        \item 【自适应步长策略】
        \item 【结合二阶信息的加速方案】
        \item 【推广到更一般的问题类】
    \end{enumerate}
\end{itemize}

\subsection{与其他方法的比较}

简要对比本文方法与文献中的相关方法，突出本质区别（不仅是数值好坏的比较，更要讨论\emph{数学结构上的差异}）。

\begin{table}[htbp]
    \centering
    \caption{与相关方法的比较}
    \label{tab:comparison}
    \begin{tabular}{lccc}
        \toprule
        方法 & 核心数学工具 & 优势 & 局限 \\
        \midrule
        本文方法      & --- & --- & --- \\
        文献 A        & --- & --- & --- \\
        文献 B        & --- & --- & --- \\
        \bottomrule
    \end{tabular}
\end{table}

\subsection{与课程内容的呼应}

反思本文所用数学方法在课程知识体系中的位置：

\begin{itemize}
    \item 本文的数学工具是课程中哪些主题的自然延伸？
    \item 课程中的哪些内容帮助了你理解该问题？
    \item 课程讲述的方法在本问题上有何局限？（是否存在课程方法直接套用失败的情况？你是如何调整的？）
\end{itemize}
        % 六、个人见解
% ============================================================
% 致谢 (Acknowledgments)
% ============================================================
\section*{致谢}

% === AI 辅助声明（必须保留）===

\noindent\textbf{AI 辅助声明}（根据课程要求，必须如实说明 AI 工具的使用范围）：

\begin{itemize}
    \item 写作语言润色：【是/否，使用了何种工具】
    \item 数学推导辅助：【是/否，使用了何种工具，具体在哪些步骤】
    \item 代码生成/补全：【是/否，使用了何种工具，具体涉及哪些模块】
    \item 图表生成：【是/否，使用了何种工具】
    \item 文献检索：【是/否，使用了何种工具】
\end{itemize}

\textbf{声明：}本文的形式化推导、直观解释、实验设计和结果分析均由作者独立完成。
【如使用了 AI 工具，在此处补充具体说明，例如：\textit{"写作过程中使用了 Claude/ChatGPT 辅助语言润色；核心算法的 Python 实现中约 30\% 的代码由 AI 生成骨架后手工修改"}】

\vspace{1em}

% === 传统致谢（可选）===
% 感谢【导师/同学】在【具体方面】给予的帮助和建议。
% 感谢课程教学团队的教学与指导。
 % 致谢

% ==================== 参考文献 ====================
\bibliographystyle{ieeenat_fullname}
\bibliography{main}

% ==================== 附录（可选）====================
% \appendix
% \input{sec/X_appendix}

\end{document}
